\cotefile{157-5}
\scanpage{52}
\mspage{39}
% archive sheet number: 51

3. \uline{Catégories de modèles à \unsure{cofibrations|cofibration}}

Soit $M$ une catégorie munie de deux \unsure{ens} de flèches
\[
 W,C\subset\mathrm{Fl}\,M,\unsure{.}
\]

On suppose que $M$ contient un objet initial, \unsure{et}

\deleted{\illegible{1}} $\mathrm{Cof}_{0}$ : $W$ est \unsure{modérément} saturé i.e. contient iso et si $f,g$ composables et si deux parmi $f,g,gf$ sont dans $W$, le troisième aussi

$\mathrm{Cof}_{1}$ : $C$ \deleted{\unsure{stable}}contient iso, stable par composition et par \unsure{cochangement} \unsure{des bases|de base}

$\mathrm{Cof}_{2}$ : $TC\overset{\mathrm{def}}{=}C\cap W$ est stable par cochangement de base

$\mathrm{Cof}_{3}$ : Toute \unsure{flèche|flèches} $f$ dans $M$ s'écrit
\[
 f=pi
\]
avec $i\in C,\ p\in W$.

\inserted{Ce sont les} Axiomes de K.~S. Brown \inserted{(\unsure{forme duale})}, qui suppose de plus que les objets de $M$ sont $C$-cofibrants, i.e. $\varnothing_M\unsure{\mathbin{-}|\longrightarrow} X$ est dans $C$. Je ne ferai pas cette hypothèse. La sous-catégorie $M_C$ \deleted{\unsure{de}} de $M$ formée des
